
% ════════════════════════════════════════════
% 引言（01.引言.tex）
% ════════════════════════════════════════════
\section{引言}
\subsection{问题背景}

随着以光伏发电为代表的分布式能源大规模接入配电网，其出力受天气影响的间歇性与随机性给电网的功率平衡带来了严峻挑战。微网作为聚合分布式电源、储能设备与本地负荷的小型发用电系统\cite{lasseter2002,luzongxiang2007}，能够通过本地消纳光伏电量、平抑负荷波动，显著减少对公共电网的依赖与冲击。在非孤岛运行方式下，微网既可以从外部电网购电，也可以将多余的电量反馈电网或存储于储能设备之中，其能量管理与运行调度已成为研究热点\cite{kanchev2011,wuxiong2014}。

微网与外网调控的核心矛盾在于：一方面，小区负荷与光伏出力在一天之内、一年之内都随时间剧烈变化，光伏只在白天出力且峰谷差异巨大；另一方面，外网电价存在明显的峰谷价差，波动电价下价差更为显著\cite{zhangqin2008}。若微网完全按负荷实时从外网购电，将承担高昂的峰时购电费用；而储能设备的引入，使微网可以在电价较低或光伏有剩余电量的时段充电，在电价较高或光伏不足的时段放电，从而在满足负荷供电的前提下尽可能降低购电费用。然而，储能设备的容量有限、充放电功率有限且存在效率损耗，更为关键的是，购电计划必须在当天负载与光伏实际发生之前制定，计划与兑现之间存在无法回避的时间差。本文要处理的就是这一时间差：当只有预测信息可用时，如何安排逐日、逐时间段的购电与充放电计划，使全年购电支出尽可能低，并在预测失准时把缺电损失限制在可控范围内。

\subsection{国内外研究现状}

微网概念与运行框架由 Lasseter 等提出\cite{lasseter2002}，此后围绕微网能量管理形成了大量研究。鲁宗相与王彩霞等系统梳理了微电网的组成、控制与保护问题\cite{luzongxiang2007}，吴雄等则从能量管理系统的角度总结了日前调度、实时调度与多时间尺度协调的实现方式\cite{wuxiong2014}。Kanchev 等针对以光伏为主电源的微网，将功率预测与运行计划结合，验证了预测信息对降低运行费用的作用\cite{kanchev2011}。这类研究确立了"预测—调度"的基本范式，但多数算例将调度周期压缩在数日之内，对整个年度的逐日滚动与跨日储电量传递讨论较少。

储能是微网消纳光伏、赚取峰谷价差的关键环节。周林等综述了适用于微网的各类储能技术及其容量配置方法\cite{zhoulin2011}，桑丙玉等以平抑新能源出力波动为目标给出了储能优化配置模型\cite{sangbingyu2014}，Sioshansi 等则在电力市场背景下估算了储能在电价套利中的价值与经济边界\cite{sioshansi2010}。这些工作表明储能收益主要来自峰谷价差，但也普遍假定储能可以自由调度，未把"计划刚性执行、实际偏差按数倍电价结算"这类结算规则纳入模型，因而难以回答计划偏离本身的经济代价。

预测是本文模型的主要输入。Antonanzas 等对光伏功率预报的方法体系与误差来源作了系统评述，指出小时级预报的误差在不同天气类型下差异明显\cite{antonanzas2016}；刘卫亮等考虑雾霾因素建立了光伏发电超短期功率预测模型\cite{liuweiliang2018}。预报误差的统计特性直接决定随机优化的情景构造，本文正是借助历史预报残差来给出误差的经验分布。

在优化方法层面，Dantzig 建立的线性规划理论\cite{dantzig1963}与 Boyd 等总结的凸优化方法\cite{boyd2004}为本文提供了建模与求解基础；Birge 与 Louveaux 系统阐述了两阶段随机规划的结构与求解途径\cite{birge2011}，司守奎等给出了常用的建模与算法实现\cite{si2011}。Parisio 等将模型预测控制用于微网经济运行\cite{parisio2014}，Rawlings 与 Camacho 等则完善了滚动时域优化的理论框架\cite{rawlings2017,camacho2013}，张钦等从需求响应角度分析了电价对用户用电行为的影响\cite{zhangqin2008}。上述方法为本文提供了直接工具，但已有微网算例多假定预测误差分布已知，或把多时段决策一次性求解，对"计划允许修改、修改按分段价格收费"这一结构缺乏显式的线性化处理。

综合来看，既有研究在微网能量管理与光伏预测方面成果丰富，但仍有两点不足：其一，部分模型在日前优化中直接使用当天的实际序列，隐含了事后信息，据此得到的计划在真实运行时难以复现；其二，对预测误差只关注其期望影响，较少把缺电损失与计划调整代价写进同一个费用函数中加以权衡。本文以一个小区微网的全年运行数据为对象，沿"单日确定性—全年随机—滚动调整—波动电价"的顺序建立四个递进的模型，试图补上这一环。

\subsection{问题重述}

题目给出一个含光伏发电、储能设备与外部电网的小区微网，附件提供了 10 分钟粒度的负荷与光伏出力序列、分时电价与波动电价数据，以及光伏出力的小时级预报。四问的共同目标是在满足负荷供电与储能运行约束的前提下给出购电与充放电计划，使购电费用尽可能低，区别在于每个时点可用的信息量与偏差的结算规则依次放宽。

\textbf{问题一：}已知全天的负荷、光伏出力与分时电价，待求各 10 分钟区间的购电功率与储能充放电功率，并要求储电量在 0:00 与 24:00 相等。该问不含任何不确定性，是后续模型的基准情形。

\textbf{问题二：}已知全年分时电价以及 10 分钟粒度的负荷、光伏序列，但在每天 0:00 制定计划时这些实际值尚未发生，只能依据预测给出计划，若实际供电不足则按 5 倍电价紧急购电。待求全年逐日的计划购电量、充放电计划及其费用。难点在于计划必须与当天真值无关，而费用由真值结算，两者之间的缺口直接决定紧急购电费量级。

\textbf{问题三：}在问题二的基础上，可于 0:00、6:00、12:00、18:00 获得未来 24 小时的小时级光伏预报。待求含计划购电、调整购电、违约惩罚与紧急购电的完整购电策略，并需判断增设其他预报时刻是否有必要。难点在于把小时级预报展开为 10 分钟序列，并用合适的惩罚结构刻画计划偏离的经济代价。

\textbf{问题四：}外网电价由分时电价改为逐 10 分钟的实时波动电价，其余条件不变。待求波动电价下当天的购电策略，并按新电价重算问题二与问题三，以考察电价波动对费用结构与调整收益的影响。

\subsection{本文技术路线}

针对上述四问，本文按"由确定到不确定、由单日到全年、由固定电价到波动电价"的顺序建立一套相互衔接的模型。问题一以 10 分钟为步长把一天离散为 144 个区间，用线性规划求单日最优购电与储能充放电计划，并标定储能参数的可用区间；问题二把单日模型扩展为"预测—决策—结算"三段式，先由扩张窗口逐日标定负荷与光伏预测模型，再以 0:00 可得的预测值与历史残差构造情景，建立两阶段随机线性规划，运行时储能按计划刚性执行，实际偏差以 5 倍电价结算；问题三在问题二的框架上加入四个预报时刻的滚动重优化，用上境变量把"计划下调 0.5 倍、上浮 1.5 倍"的分段惩罚精确线性化；问题四把分时电价替换为实时波动电价，重算问题二与问题三，比较电价波动对费用结构的影响。四问共用同一套数据预处理流程与储能参数设置，模型之间的关系与求解步骤如图~\ref{fig:overall_analysis}~所示。
